\documentclass[11pt]{article}

\usepackage[margin=1.1in]{geometry}
\usepackage{amsmath,amssymb}
\usepackage{booktabs}
\usepackage{microtype}
\usepackage[colorlinks=true,linkcolor=blue,citecolor=blue,urlcolor=blue]{hyperref}

\newcommand{\Itp}{\mathit{Itp}}
\newcommand{\true}{\top}
\newcommand{\false}{\bot}
\newcommand{\lbla}{\mathsf{a}}
\newcommand{\lblb}{\mathsf{b}}
\newcommand{\lblab}{\mathsf{ab}}
\newcommand{\opt}[1]{\texttt{#1}}
\newcommand{\code}[1]{\texttt{#1}}

\title{EUF Craig Interpolation in Z3 via DRUP Proof Replay:\\
Reordering, Labeled Systems, and Colorability}
\author{Yakir Vizel \and Claude (Anthropic)}
\date{July 2026}

\begin{document}
\maketitle

\begin{abstract}
This note documents the Craig interpolation pipeline implemented in a
development branch of Z3 for quantifier-free EUF. Interpolants are computed
by replaying a trimmed clausal (DRUP-style) proof, combining a labeled
propositional interpolation system with a theory-specific procedure for EUF
lemmas based on proof-forest summarization. On top of the basic pipeline the
branch implements the proof reordering techniques of \emph{DRUPing for
Interpolants} (colored unit propagation and restructuring of resolution
chains into single-colored stages), the labeled interpolation systems of
D'Silva et al.\ (McMillan, HKP, and McMillan's dual, with a differential
strength check), and a per-variable label optimization that targets
colorability of theory lemmas. We describe how the pieces fit together,
what was proved and measured along the way, and directions for future work.
\end{abstract}

\section{Introduction}

Given an unsatisfiable pair of formulas $(A, B)$, a Craig interpolant is a
formula $I$ over the symbols common to $A$ and $B$ such that $A \models I$
and $I \wedge B$ is unsatisfiable. The implementation described here
computes interpolants for quantifier-free EUF problems from the clausal
proofs produced by Z3's \code{sat.euf} solver, following the general recipe
of \emph{DRUPing for Interpolants}~\cite{drup-itp}: log a DRUP-style proof
during solving, trim it, and compute the interpolant during a validated
replay of the trimmed proof. Theory reasoning enters the proof as EUF lemmas
with certificates (hints), which receive their own partial interpolants from
a dedicated EUF interpolation procedure.

The relevant solver options are collected in
Table~\ref{tab:options}; the directory \code{euf\_examples/} in the
repository contains the instances referenced throughout, each of which
passes the built-in interpolant validity checks.

\section{The Interpolation Pipeline}
\label{sec:pipeline}

\subsection{Proof logging and partitions}

The user partitions the input with \code{(set-itp-group A)} and
\code{(set-itp-group B)}; with \opt{solver.proof.interpolate\_log} enabled,
every logged input clause carries its partition as a marker
(\code{\_\_itp\_A} / \code{\_\_itp\_B}). During solving, the clause proof is
logged (\opt{solver.proof.log}): input clauses as \code{assume} steps, RUP
inferences as \code{infer \ldots\ rup}, EUF lemmas as inferences carrying an
\code{euf} hint that certifies the lemma (congruence and transitivity
premises), and clause deletions as \code{del} steps.

\subsection{Trimming and marking}

The proof is replayed by \code{proof\_trim}
(\opt{solver.proof.trim}), which re-runs the proof steps forward and then
performs the backward pass of~\cite{drup-itp}: starting from the terminal
conflict, each in-core inference is justified by a conflict-analysis probe
(assert the negated clause, propagate, trace the conflict cone), which marks
the clauses it depends on as core and records them as the inference's
dependencies. Theory lemmas are treated as \emph{assumptions} during
trimming --- they are not RUP --- with their full clause reconstructed by the
proof checker when base-level-false literals were dropped from the logged
clause.

When interpolation is enabled (\opt{solver.proof.interpolate}), trimming
additionally computes the $A$/$B$ \emph{markings} used by everything
downstream. A marking is a subset of $\{A, B\}$, combined by union:
\begin{itemize}
\item an original clause is marked $A$ or $B$ by its partition;
\item an inferred clause is marked with the union of the marks of its
      antecedents;
\item a variable is marked with the union of the marks of the original core
      clauses in which it occurs (mark $AB$ = \emph{shared});
\item a level-0 trail literal is marked by its unit derivation.
\end{itemize}

\subsection{Replay and validation}

The trimmed proof is replayed by \code{proof\_replay\_validator} on a fresh
internal SAT solver. Assumptions are added directly; every RUP inference is
validated by a conflict-analysis probe, and the resulting \emph{trivial
resolution chain} --- the sequence of reason clauses and pivots that derives
the inference --- is handed to a visitor. Unit propagations at level~0 are
likewise traced and emitted as chains. The visitor abstraction decouples
replay from interpolation: the same replay drives logging, validation, and
interpolant construction.

\subsection{Interpolant construction}

The interpolating visitor (\code{itp\_visitor}) maintains a partial
interpolant for every derived clause and unit. Input clauses receive leaf
labels, theory lemmas receive EUF partial interpolants
(Section~\ref{sec:euf-leaf}), and every resolution chain combines the
partial interpolants of its premises pivot by pivot according to the labeled
system of Section~\ref{sec:labelings}. The partial interpolant of the empty
clause is the interpolant; it is simplified and, with
\opt{solver.proof.check\_interpolant}, verified by three solver checks:
$A \Rightarrow I$, $I \wedge B$ unsatisfiable, and that $I$ uses only shared
symbols.

\subsection{Theory lemmas: EUF partial interpolants}
\label{sec:euf-leaf}

An EUF lemma $C$ is a valid clause; its negation $\neg C$ is an
unsatisfiable conjunction of equalities and disequalities. The literals of
$\neg C$ are charged to an $(\alpha, \beta, \gamma)$ triple according to
their labels (Section~\ref{sec:labelings}); the lemma's partial interpolant
$J$ must satisfy the \emph{labeled T-lemma obligations}
\begin{equation}
\label{eq:tlemma}
\alpha \wedge \gamma \models_{\mathrm{EUF}} J
\qquad\text{and}\qquad
J \wedge \beta \wedge \gamma \models_{\mathrm{EUF}} \false ,
\end{equation}
over shared symbols. $J$ is computed by \code{euf\_interpolator} using
colored proof-forest summarization:
\begin{enumerate}
\item build an egraph over all terms of the lemma; merge the equalities of
      the \emph{summarized} side first, with their proof-forest edges
      \emph{marked}, and close under congruence; then merge the other
      side's equalities unmarked;
\item locate a violated disequality and walk the proof-forest path between
      its endpoints (\code{euf\_summarizer}): every maximal run of marked
      edges is summarized by an equality between its endpoint terms;
      congruence edges are re-derived recursively, folded into runs when
      their function symbol and argument paths are usable in the
      summarized side, and otherwise replaced by their argument equalities
      (introducing shared boundary terms);
\item a conflict on a $\beta$ disequality is interpolated by summarizing
      the $\alpha\gamma$ side; a conflict on an $\alpha$ disequality
      dually, with $J$ negated ($\Itp(A,B) = \neg \Itp(B,A)$).
\end{enumerate}
Trivial cases short-circuit: if $\beta \wedge \gamma$ is unsatisfiable
alone, $J = \true$; if $\alpha \wedge \gamma$ is, $J = \false$.

\subsection{Stitching}

In summary, the propositional skeleton (resolution over input clauses,
theory lemmas, and derived clauses) is interpolated by the labeled system,
while each theory lemma contributes a leaf interpolant satisfying
\eqref{eq:tlemma}. Soundness of the composition rests on the leaf and pivot
rules maintaining, for every derived clause $C$ with partial interpolant
$J(C)$, the invariant
$A \models C{\downarrow}\lbla \vee J(C)$ and
$B \models C{\downarrow}\lblb \vee \neg J(C)$,
where $C{\downarrow}\ell$ restricts $C$ to its $\ell$-labeled literals; the
theory obligations \eqref{eq:tlemma} are exactly this invariant for T-valid
leaves.

\section{Proof Reordering and Colorability}
\label{sec:reorder}

A resolution proof is \emph{colorable} when every chain resolves clauses of
a single partition; colorable proofs yield structurally simple (CNF-like)
interpolants, whereas chains that interleave $A$ and $B$ steps nest the
$\vee/\wedge$ combinations of the labeling one level per color crossing.
Following~\cite{drup-itp}, the replay reorders proofs toward colorability in
two places, both gated by \opt{solver.proof.reorder} (default on).

\subsection{Colored unit propagation}

The replay solver's propagation runs in \emph{rounds} driven by a clause
filter: round~1 propagates only $A$-marked clauses to fixpoint, round~2
admits every clause; each round restarts the queue scan from the same
initial position, and a clause skipped in an early round keeps its watch
entries untouched. Because the last round is unrestricted, filtered
propagation assigns exactly the same literals as plain propagation --- only
the \emph{order}, and hence the reasons recorded on the trail, change. The
trail itself is never permuted, so assignment order remains
antecedent-before-consequent, the invariant on which replay, tracing, and
backtracking rely.\footnote{An earlier implementation reordered the trail in
place; it corrupted both the interpolant construction (out-of-order unit
tracing) and backtracking (literals stranded across scope boundaries), and
was replaced by the round-based scheme. A related, independent defect was
found and fixed in the replay indexing: proof step identifiers do not
coincide with trail positions once deletions or duplicate units occur, and
both replay paths now map identifiers to trail positions explicitly.}
An orthogonal option, \opt{solver.proof.core\_first\_bcp}, adds leading
rounds restricted to clauses already marked core, biasing new chains toward
reusing the existing core, with an unrestricted final round preserving
completeness.

\subsection{Restructuring RUP chains}

Chain construction is \emph{partition-bounded} conflict analysis: clauses
are ranked $\mathrm{part}(C) \in \{1, 2\}$ (pure-$A$ versus everything
else), and each lemma is derived in stages. If the conflict lies within
partition~1, a first pass resolves only pivots with partition-1 reasons; if
it already derives the lemma, the chain is single-colored and we are done,
and otherwise the result is emitted as an \emph{intermediate lemma} from
which the analysis restarts with the bound lifted. Inside a mixed chain, a
pivot whose reason is pure-$A$ first has that reason \emph{purified}
(MiniSat's \code{fixrec}): it is re-derived within partition~1 and emitted
as its own single-colored chain, and the mixed chain uses the purified
clause as a leaf. Purified reasons are cached per variable assignment.
Pivots are always resolved in reverse trail order, so chains remain trivial
resolution chains and no pivot is ever re-introduced.

\subsection{Colorability on the EUF side}

The natural question is whether the same reordering applies inside theory
lemmas. The answer, somewhat surprisingly, is that the EUF procedure of
Section~\ref{sec:euf-leaf} \emph{already} contains both halves of the
recipe, in a stronger form:

\begin{itemize}
\item \textbf{Merge order is colored propagation.} Closing the summarized
      side first, with marked justifications, makes the proof forest
      \emph{maximally} colored for that side: any equality it entails is
      derived --- and marked --- in the first phase, because the acyclic
      union-find never re-connects an already-connected pair. Congruence
      closure is complete for ground EUF, so this single eager closure
      captures everything at once, where propositional colored BCP must
      still restructure chains after the fact.
\item \textbf{Congruence re-derivation is purification.} Congruence edges
      are created unmarked; the summarizer re-derives each one, marks it
      when its symbol and argument paths are usable in the summarized side
      (fusing marked runs --- exactly what purification would do), and
      otherwise replaces it by argument equalities, which is the best
      colorable summary available.
\end{itemize}

Consequently a path-level purification pass --- re-marking runs of
other-side edges whose endpoints the summarized side entails on its own ---
is provably a no-op under the current merge order. The pass is nonetheless
implemented, guarded by a side-entailment oracle, as an \emph{executable
statement of the invariant}: if a future change (merge order, justification
marking, non-minimal lemma explanations) breaks the argument, purification
becomes live and reports itself instead of silently degrading interpolants.
Measurements confirm zero activations across the example suite, including
with the merge phases deliberately reversed (trimmed lemmas are minimal
conflicts and admit no redundant justification choices).

\subsection{Effect}

On \code{itp\_reorder.smt2} --- two thirty-step implication chains over
equality atoms with color pattern $A\,A\,B$ repeating --- the raw
interpolant with reordering is a flat, CNF-like conjunction with one small
clause per color-run boundary; with \opt{solver.proof.reorder=false} it is a
deeply nested $\vee/\wedge$ alternation of the same size (179 nodes) that
simplification cannot flatten. On tree-shaped EUF ladders the final
interpolants coincide with and without reordering: $\wedge/\vee$ are
associative--commutative, so re-association alone vanishes under
simplification; the benefit there is structural (single-colored chains,
factored $A$-lemmas).

\section{Labeled Interpolation Systems}
\label{sec:labelings}

The propositional construction is a labeled interpolation system in the
sense of D'Silva et al.~\cite{lis}. Every literal occurrence carries a label
in the lattice $\lblb \sqsubseteq \lblab \sqsubseteq \lbla$; locality forces
$A$-local variables to $\lbla$ and $B$-local variables to $\lblb$, while
shared variables are free. The implementation uses uniform per-variable
labels: a base policy (\opt{solver.proof.itp\_labeling}) fixes the label of
shared variables ---
\begin{center}
\begin{tabular}{lll}
\toprule
policy & shared label & system \\
\midrule
\opt{mcmillan} (default) & $\lblb$ & McMillan~\cite{mcmillan} (strongest) \\
\opt{hkp}                & $\lblab$ & Huang / Kraj\'{\i}\v{c}ek / Pudl\'ak (symmetric) \\
\opt{dual}               & $\lbla$ & McMillan's dual (weakest) \\
\bottomrule
\end{tabular}
\end{center}
--- and unmarked variables are treated as $\lblb$. The rules are:
\begin{align*}
C \in A:&\quad J(C) = \bigvee \{\, \ell \in C : L(\ell) = \lblb \,\} \\
C \in B:&\quad J(C) = \bigwedge \{\, \neg\ell : \ell \in C,\ L(\ell) = \lbla \,\} \\
\text{pivot } p:&\quad
\begin{cases}
J_1 \vee J_2 & L(p) = \lbla \\
J_1 \wedge J_2 & L(p) = \lblb \\
(J_1 \vee \neg p) \wedge (J_2 \vee p) & L(p) = \lblab
\end{cases}
\end{align*}
where, in the $\lblab$ rule, the premise of $J_2$ contains the pivot literal
$p$. Under the default policy this coincides literally with McMillan's
system, and the implementation is bit-for-bit backward compatible.

\paragraph{Theory lemmas under labels.}
The labels determine the $(\alpha, \beta, \gamma)$ split of
Section~\ref{sec:euf-leaf}: $\lbla$-labeled atoms go to $\alpha$,
$\lblb$-labeled to $\beta$, and $\lblab$-labeled to $\gamma$, which is
available to \emph{both} sides per the obligations \eqref{eq:tlemma}.
Operationally, $\gamma$ atoms join whichever side is summarized (merged
marked in the first phase, and included in the purification oracle), which
is also the colorability-optimal placement. A subtle point uncovered during
this work: charging a shared atom to $\alpha$ while resolving its pivot with
the $\lblb$ rule (as an earlier version did) is a mixed labeling combined
with a pure-$\lblb$ pivot rule and is unsound in general; the split must
follow the labels, which the current implementation enforces by
construction.

\paragraph{Strength ordering.}
For a fixed proof, pointwise-stronger labelings yield logically stronger
interpolants; in particular
$\Itp_{\text{mcmillan}} \models \Itp_{\text{hkp}} \models
\Itp_{\text{dual}}$. The option \opt{solver.proof.check\_labeling\_order}
recomputes the interpolant under all three labelings on the same trimmed
proof and verifies both implications with solver calls. Theory-lemma
partial interpolants are not strength-ordered leaf-wise across labelings,
so with theory lemmas present a violation is reported as advisory; on the
current suite the order holds strictly everywhere, and on
\code{itp\_reorder.smt2} both reverse implications fail, i.e.\ the three
systems are strictly separated (179 / 209 / 179 raw nodes with three
distinct shapes).

\section{Colorability-Directed Labeling}
\label{sec:labelopt}

The $\lblab$ label is a double-edged sword: inside a theory lemma it lets a
shared equality join the summarized side, keeping the lemma's internal chain
single-colored; at every propositional resolution on that variable it costs
an $\lblab$-guard $(J_1 \vee \neg p) \wedge (J_2 \vee p)$. Uniform policies
take both effects or neither. The per-variable optimization
\opt{solver.proof.itp\_label\_opt=colorable} takes only the good half: it
promotes to $\lblab$ \emph{exactly the shared variables that occur in theory
lemmas}, leaving all other shared variables at the base policy's label.
Labels are computed once, before replay, so every occurrence of a variable
is labeled consistently (which soundness requires), and promotion is
monotone in the label lattice, so the strength order check of
Section~\ref{sec:labelings} remains valid.

The example \code{colorable\_opt.smt2} combines a congruence chain through a
shared link with a colored propositional implication chain in a single
conflict and separates the configurations:

\begin{center}
\begin{tabular}{lcl}
\toprule
configuration & raw nodes & lemma interpolant \\
\midrule
\opt{mcmillan} & 37 & splits at the shared boundary \\
\opt{mcmillan} + \opt{colorable} & 39 & single colorable equality \\
\opt{hkp} & 47 & single equality, guards on every chain boundary \\
\bottomrule
\end{tabular}
\end{center}

The optimization degenerates correctly at both extremes: it coincides with
\opt{hkp} when every shared variable occurs in a lemma
(\code{theory\_ab2.smt2}) and with \opt{mcmillan} when none does
(\code{itp\_reorder.smt2}).

\section{Future Work}
\label{sec:future}

\begin{itemize}
\item \textbf{Size- and strength-directed labeling.} The label space is
      larger than one heuristic: replays on a trimmed proof are cheap, so a
      search over per-variable promotions (keep a promotion if the raw
      interpolant shrinks) is practical. For model-checking clients,
      strength can matter more than size; exposing the labeling per query
      would let e.g.\ an IC3/Avy-style loop pick weak or strong
      interpolants adaptively.
\item \textbf{Per-occurrence labels.} The unit-label cache is compatible
      with per-occurrence labels (partial interpolants attach to
      derivations, not occurrences); what is missing is a policy surface
      and a use case, e.g.\ PeRIPLO-style proof-sensitive labeling.
\item \textbf{Sequence interpolants.} The visitor design mirrors Avy's
      \code{ItpSequence}; generalizing from one $(A,B)$ cut to a sequence
      of cuts over the same replay would produce sequence interpolants in
      one pass, which is what BMC-based model checkers consume.
\item \textbf{Richer theory coverage.} \code{euf\_interpolator} handles
      equality atoms; uninterpreted predicates and other convex theories
      (e.g.\ difference logic) could reuse the same
      $(\alpha,\beta,\gamma)$ leaf interface. Non-equality atoms currently
      fall back to a generic leaf, which weakens interpolants; the
      fallback should at least be instrumented.
\item \textbf{Interpolant post-processing.} Summaries can contain
      duplicate or subsumed conjuncts across chains; a dedicated
      simplification (or emitting summaries into a hash-consed set) would
      shrink outputs beyond what generic rewriting achieves.
\item \textbf{Scaling studies.} The example suite is synthetic and small.
      Running the pipeline on Avy/HWMCC-derived IUC queries would exercise
      colored propagation and chain restructuring where they matter
      (large, deletion-heavy proofs) and quantify the effect of
      \opt{core\_first\_bcp} and \opt{colorable} on interpolant size and
      convergence.
\item \textbf{Formalization.} The labeled T-lemma obligations
      \eqref{eq:tlemma}, the maximality argument for the two-phase merge,
      and the soundness of chain restructuring are proved on paper (in
      code comments); a mechanized or at least carefully written-up proof
      of the composed system would be valuable, particularly around the
      interaction of purification, labels, and the dual (negated) EUF
      case.
\end{itemize}

\appendix

\section{Option summary}

\begin{table}[h]
\centering
\small
\begin{tabular}{lp{8.2cm}}
\toprule
option & effect \\
\midrule
\opt{solver.proof.log} & log the clause proof to a file \\
\opt{solver.proof.interpolate\_log} & emit \code{\_\_itp\_A}/\code{\_\_itp\_B} partition markers \\
\opt{solver.proof.trim} & trim the logged proof \\
\opt{solver.proof.replay} & replay with validation after trimming \\
\opt{solver.proof.interpolate} & compute markings and an interpolant during replay \\
\opt{solver.proof.check\_interpolant} & verify $A \Rightarrow I$, $I \wedge B$ unsat, shared symbols \\
\opt{solver.proof.reorder} & colored propagation + chain restructuring (default on) \\
\opt{solver.proof.core\_first\_bcp} & prefer core-marked clauses in replay propagation \\
\opt{solver.proof.itp\_labeling} & \opt{mcmillan} / \opt{hkp} / \opt{dual} shared-variable labeling \\
\opt{solver.proof.itp\_label\_opt} & \opt{none} / \opt{colorable} per-variable label optimization \\
\opt{solver.proof.check\_labeling\_order} & verify $\Itp_{\text{mc}} \models \Itp_{\text{hkp}} \models \Itp_{\text{dual}}$ \\
\bottomrule
\end{tabular}
\caption{Interpolation-related options.}
\label{tab:options}
\end{table}

\begin{thebibliography}{9}

\bibitem{drup-itp}
A.~Gurfinkel and Y.~Vizel.
\newblock DRUPing for interpolants.
\newblock In \emph{Formal Methods in Computer-Aided Design (FMCAD)}, 2014.

\bibitem{lis}
V.~D'Silva, D.~Kroening, M.~Purandare, and G.~Weissenbacher.
\newblock Interpolant strength.
\newblock In \emph{Verification, Model Checking, and Abstract Interpretation
  (VMCAI)}, 2010.

\bibitem{mcmillan}
K.~L. McMillan.
\newblock Interpolation and SAT-based model checking.
\newblock In \emph{Computer Aided Verification (CAV)}, 2003.

\end{thebibliography}

\end{document}
